\section*{Chapter \ref{ch:b3:scale-variety-firms} --- Scale, Variety and Firms}

\begin{solution}{exo:b3:scale-variety-firms:1}
Largest: computers and electronics and electrical equipment, 38~\%; smallest:
miscellaneous manufacturing, 2~\%. % figdata: exporters
Above 18~\%: computers and electronics, electrical equipment, chemicals (36),
machinery (33), primary metals (30), plastics and rubber and transport equipment
(28), textile mills (25), leather and paper (24), beverages and tobacco (23);
petroleum and coal equals it. Most are capital- or skill-intensive industries,
in which a skill-abundant country such as the United States has a \omterm{def:b3:ricardian-trade:comparative-advantage}{comparative
advantage}: endowments still matter for which industries export most.
\end{solution}

\begin{solution}{exo:b3:scale-variety-firms:2}
Furniture $1 - 10/70 = 0.857$; machines $1 - 40/60 = 0.333$; grain $1 - 60/60 = 0$.
Aggregate: $1 - (10 + 40 + 60)/(70 + 60 + 60) = 1 - 110/190 = 0.421$. Ricardo's model
explains the one-way trade in grain; the two-way trade in furniture needs
varieties.
\end{solution}

\begin{solution}{exo:b3:scale-variety-firms:3}
$p = \sigma c/(\sigma - 1) = \tfrac43 c$; Lerner index $(p - c)/p = 1/\sigma = 0.25$. Closer
substitutes mean a higher $\sigma$: the markup and the Lerner index fall toward 1 and 0,
the competitive values.
\end{solution}

\begin{solution}{exo:b3:scale-variety-firms:4}
A country with a larger home market for a differentiated, costly-to-ship good
that is a net importer of it --- or whose share of the industry's production is
smaller than its share of demand --- contradicts it. In the \omterm{def:b3:factor-endowments:ho}{Heckscher--Ohlin model},
with production set by endowments and prices equalised by trade, a country whose
consumers like a good more consumes more of it than it makes and imports it; in
the Krugman model, demand draws firms in, and the country exports it.
\end{solution}

\begin{solution}{exo:b3:scale-variety-firms:5}
$p = 5/4 = 1.25$; $\bar x = (\sigma - 1)F/c = 40$; employment per firm $\sigma F = 50$;
$n = 1{,}000/50 = 20$. Real wage $1/P = n^{1/(\sigma-1)}/p = 20^{1/4}/1.25 = 2.115/1.25 =
1.692$.
\end{solution}

\begin{solution}{exo:b3:scale-variety-firms:6}
40 varieties, 20 from each country; price still 1.25. Real wage $40^{1/4}/1.25$, a
factor $2^{1/4} = 1.189$: an 18.9~\% gain. With $\sigma = 3$ the factor is $2^{1/2} =
1.414$, a 41.4~\% gain: when varieties are less alike, each new one is worth more.
\end{solution}

\begin{solution}{exo:b3:scale-variety-firms:7}
The tested solution is \code{ex07_cutoff.py}, in the chapter's \code{exercises/}
folder. The formula gives $\varphi^{*6} = 10 \times 2/(4 \times 1) = 5$, $\varphi^* = 5^{1/6} =
1.3077$, and the bisection on a million draws gives 1.3065. A share $1/5 = 20~\%$ of
entrants produce (20.1~\% in the sample). % output: ex07_cutoff
\end{solution}

\begin{solution}{exo:b3:scale-variety-firms:8}
Entrants that produce: $1.5^{-3} = 0.296$. Producers that export: $\Pr(\varphi > 3 \mid
\varphi > 1.5) = (1.5/3)^3 = 0.125$. With $k = 2$, $(1/2)^2 = 0.25$: a thicker tail of
productive firms means more exporters.
\end{solution}

\begin{solution}{exo:b3:scale-variety-firms:9}
Follow firms over time and compare the productivity of future exporters, in the
years before they start exporting, with that of firms that never export, in the
same industry. If future exporters were no more productive before they started,
and became so only after, the data would contradict selection and support
learning by exporting: firms gaining from contact with foreign buyers and
technology.
\end{solution}

\begin{solution}{exo:b3:scale-variety-firms:10}
Write $\rho = (\sigma - 1)/\sigma$ and maximise $\sum_i x_i^\rho$ subject to $\sum_i p_i x_i = E$.
The first-order conditions $x_i^{\rho - 1} = \mu p_i$ give $x_i = (\mu p_i)^{-\sigma}$, so
$x_i/x_j = (p_i/p_j)^{-\sigma}$. Substituting into the budget, $\mu^{-\sigma}\sum_j p_j^{1-\sigma} = E$,
so $\mu^{-\sigma} = E/P^{1-\sigma}$ and $x_i = p_i^{-\sigma}E/P^{1-\sigma}$. Then $\sum_i x_i^\rho =
(E/P^{1-\sigma})^\rho \sum_i p_i^{1-\sigma} = E^\rho P^{(1-\sigma)(1-\rho)}$, and since $(1-\sigma)(1-\rho) =
(1-\sigma)/\sigma = -\rho$, $U = \bigl(E^\rho P^{-\rho}\bigr)^{1/\rho} = E/P$. With $P$ given, demand has
constant elasticity $\sigma$, and the monopoly rule $(p - c)/p = 1/\sigma$ gives the
markup.
\end{solution}

\begin{solution}{exo:b3:scale-variety-firms:11}
The tested solution is \code{ex11_industry_productivity.py}, in the chapter's
\code{exercises/} folder. Cut-offs: 1.246 closed, 1.293 open, 1.924 for exports.
Industry productivity rises from 2.169 to 2.387, by 10.0~\%. Among the survivors
alone, autarky productivity is 2.240: exit accounts for 33~\% of the rise, and the
other 67~\% is market share moving to exporters. % output: ex11_industry_productivity
\end{solution}

\begin{solution}{exo:b3:scale-variety-firms:12}
The term $\tau^{-4}f/f_x$ rises from $0.4823 \times 0.4444 = 0.2143$ to $0.6830 \times 0.4444
= 0.3036$, so $\varphi^{*4} = 16 \times 1.3036 = 20.86$ and $\varphi^* = 2.137$. The \omterm{def:b3:scale-variety-firms:melitz}{export cut-off}
is $1.1 \times 1.5 \times 2.137 = 3.526$, and the share of producers that export rises from
9.5~\% to $(2.137/3.526)^4 = 13.5~\%$. The firm with $\varphi = 4$ exported $120 \times
(4/2.0995)^2/1.44 = 302.5$; now $120 \times (4/2.137)^2/1.21 = 347.4$, 14.9~\% more. Both
margins rise. When $f_x$ is cut instead (problem, question 18), the extensive
margin rises but each existing exporter sells less abroad: the fixed cost does
not enter the price, while the tougher competition from new exporters lowers
everyone's sales.
\end{solution}

\begin{solution}{pb:b3:scale-variety-firms:furniture}
\textbf{1.} $p = \tfrac{3}{2}\cdot\tfrac{1}{\varphi}$: a markup of 1.5. % figdata: furn_params, furn_results

\textbf{2.} Profit $r/\sigma - f = 0$, so $r = \sigma f = 120$.

\textbf{3.} $r(\varphi) = 120(\varphi/\varphi^*)^2$, so profit $r/3 - 40 = 40\bigl((\varphi/\varphi^*)^2 - 1\bigr)$.

\textbf{4.} $40\varphi^{*-4} \times 2/2 = 2.5$, so $\varphi^{*4} = 16$ and $\varphi^* = 2$.

\textbf{5.} $2^{-4} = 1/16 = 6.25~\%$. Average productivity $\tfrac43 \times 2 = 2.667$.

\textbf{6.} Revenue $120 \times (4/2)^2 = 480$; profit $480/3 - 40 = 120$.

\textbf{7.} Exports bring $\lambda r(\varphi)/3 - 90$ with $\lambda = 1.2^{-2}$; at the cut-off
$1.2^{-2} \times 40(\varphi_x/\varphi^*)^2 = 90$, so $(\varphi_x/\varphi^*)^2 = 1.44 \times 2.25 = 3.24$ and
$\varphi_x/\varphi^* = 1.8$.

\textbf{8.} $\tau^{-4}f/f_x = 0.4823 \times 40/90 = 0.2143$; $\varphi^{*4} = 16 \times 1.2143 = 19.43$,
$\varphi^* = 2.0995$; $\varphi_x = 1.8 \times 2.0995 = 3.779$.

\textbf{9.} Those with $\varphi$ between 2 and 2.0995: $1 - 2^{-4}/2.0995^{-4} = 1 - 16/19.43 =
17.6~\%$.

\textbf{10.} $1.8^{-4} = 9.5~\%$, close to the 7~\% of United States furniture
makers in the opening chart. % ledger: b3sv:bjrs.t2.337

\textbf{11.} Domestic sales are $1.8^2 = 3.24$ times as large, and exports add a
factor $1 + 1/1.44 = 1.694$: $3.24 \times 1.694 = 5.49$.

\textbf{12.} Domestic sales $120 \times (4/2.0995)^2 = 435.6$; exports $435.6/1.44 = 302.5$;
total 738.1, against 480 in autarky, 53.8~\% more.

\textbf{13.} It still produces but does not export; its sales fall from $120 \times
(3/2)^2 = 270$ to $120 \times (3/2.0995)^2 = 245.0$, by 9.3~\%, because Kestrian varieties
now compete for Brenland's buyers.

\textbf{14.} $(\varphi_x/\varphi^*)^2 = 1.44 \times 45/40 = 1.62$, ratio 1.273. $\tau^{-4}f/f_x$ doubles to
0.4287: $\varphi^{*4} = 16 \times 1.4287 = 22.86$, $\varphi^* = 2.187$; $\varphi_x = 1.273 \times 2.187 = 2.783$.

\textbf{15.} $1.273^{-4} = 38.1~\%$.

\textbf{16.} $1 - 19.43/22.86 = 15.0~\%$.

\textbf{17.} It becomes an exporter ($3 > 2.783$). Domestic sales $120 \times (3/2.187)^2 =
225.9$, exports $225.9/1.44 = 156.9$, total 382.8, against 245.0: 56~\% more.

\textbf{18.} $120 \times (4/2.187)^2/1.44 = 278.9$, against 302.5: 7.8~\% less. The fixed
cost does not change the firm's price; but more Brenlandic and Kestrian firms now
sell in each market, and the higher cut-off is the sign of tougher competition
that lowers every incumbent's sales.

\textbf{19.} Yes. The share of producers that export quadruples (38.1~\% against
9.5~\%); producers themselves fall by 15~\%, so exporters per entrant rise 3.4 times,
from $3.779^{-4} = 0.0049$ to $2.783^{-4} = 0.0167$. Export sales per entrant are
$\sigma f \tau^{1-\sigma} (\varphi^*)^{1-\sigma}\, \tfrac{k}{k-\sigma+1}\,
\varphi_x^{\sigma-1-k}$: 2.65 thousand dalens before, 4.50 after, a rise of 70~\%. The
rise comes through the extensive margin, new exporters; the intensive margin, the
average export sales per exporter, falls from 540 to 270.

\textbf{20.} From $\tfrac43 \times 2.0995 = 2.799$ to $\tfrac43 \times 2.187 = 2.915$, by 4.1~\%
($2.1866/2.0995 = 1.041$).

\textbf{21.} From selection: the least productive firms exit and their workers
and buyers go to more productive firms; exporters' market shares grow.

\textbf{22.} Opening: $2.0995/2 = 1.050$, a 5.0~\% gain; the agreement: $1.041$, 4.1~\%;
in all $2.187/2 = 1.093$, 9.3~\%.

\textbf{23.} No change in productivity: identical firms all export, none exits,
and the gains come only from more varieties
(\cref{prop:b3:scale-variety-firms:krugman}).

\textbf{24.} Plant exit and a rise in labour productivity in the most exposed
industries, about half of it from turnover and efficiency, as the
\omterm{def:b3:scale-variety-firms:melitz}{heterogeneous-firms model} predicts. % ledger: b3sv:tre.fta

\textbf{25.} The owners of the 15~\% of firms that exit lose their entry costs and
profits, and their workers lose their jobs and must find others; in this model
they are re-employed at the same wage, while in the Canadian record employment in
the most exposed industries fell by 15~\%.

\textbf{26.} After the agreement, the average productivity of Brenland's furniture
makers rises by 4.1~\%, while 15.0~\% of the firms producing before the agreement
exit.
\end{solution}
